% TEMPLATE-MILC-v3.tex - plantilla maestra para todas las guias MILC v3
% Ver scripts/generadores/build-guia-milc.py para la lista completa de
% claves esperadas. El script valida que no queden placeholders sin
% reemplazar antes de escribir el archivo final.

\documentclass[11pt,letterpaper]{article}
\usepackage[margin=1.75cm]{geometry}
\usepackage[table]{xcolor}
\usepackage{fontspec}
\usepackage[spanish]{babel}
\usepackage{tikz}
\usepackage[most]{tcolorbox}
\usepackage{tabularx}
\usepackage{array}
\usepackage{fancyhdr}
\usepackage{titlesec}
\usepackage{ragged2e}
\usepackage{enumitem}
\usepackage{microtype}

\setmainfont{Helvetica Neue}
\setsansfont{Helvetica Neue}
\setlength{\parindent}{0pt}
\setlength{\parskip}{6pt}
\hyphenpenalty=200
\exhyphenpenalty=200
\pretolerance=400
\tolerance=2000
\setlength{\emergencystretch}{1em}
\renewcommand{\arraystretch}{1.25}
\setlist{leftmargin=5mm,itemsep=1mm,topsep=1mm}
\newcolumntype{Y}{>{\RaggedRight\arraybackslash}X}

\definecolor{milcMagenta}{HTML}{E6007E}
\definecolor{milcVino}{HTML}{5A0038}
\definecolor{milcTurquesa}{HTML}{0093A5}
\definecolor{milcVerde}{HTML}{58A923}
\definecolor{milcMostaza}{HTML}{E5B400}
\definecolor{milcOcre}{HTML}{B86600}
\definecolor{milcNegro}{HTML}{241816}
\definecolor{milcGris}{HTML}{F7F7F4}
\definecolor{milcLinea}{HTML}{E8E2DD}
\definecolor{milcGradePrimary}{HTML}{0066FF}
\definecolor{milcGradeDark}{HTML}{003D99}
\definecolor{milcGradeSoft}{HTML}{D6E8FF}
\definecolor{milcGradeAccent}{HTML}{CFE7FF}
\definecolor{milcGradeText}{HTML}{FFFFFF}
\color{milcNegro}

\pagestyle{fancy}
\fancyhf{}
\lhead{\small Tecnología e Informática · Grado 6}
\rhead{\small Guía 27 · MILC v3}
\cfoot{\small\thepage}
\renewcommand{\headrulewidth}{0pt}

\newtcolorbox{titlebox}[1]{
  enhanced,colback=#1,colframe=#1,arc=7mm,boxrule=0pt,
  left=7mm,right=7mm,top=3mm,bottom=3mm,
  before skip=8pt,after skip=8pt,
  fontupper=\sffamily\bfseries\Large\color{white}
}

\newtcolorbox{softbox}[3]{
  enhanced,colback=#2,colframe=#1,borderline west={4pt}{0pt}{#1},
  arc=2mm,boxrule=.4pt,left=4mm,right=4mm,top=3mm,bottom=3mm,
  before skip=6pt,after skip=8pt,title={#3},coltitle=white,
  colbacktitle=#1,fonttitle=\sffamily\bfseries,fontupper=\RaggedRight,
  attach boxed title to top left={xshift=3mm,yshift=-2mm},
  boxed title style={arc=2mm, boxrule=0pt}
}

\newtcolorbox{drawbox}[2]{
  enhanced,colback=white,colframe=#1,arc=4mm,boxrule=1pt,
  left=5mm,right=5mm,top=4mm,bottom=4mm,before skip=8pt,after skip=8pt,
  title={#2},coltitle=white,colbacktitle=#1,fonttitle=\sffamily\bfseries,
  fontupper=\RaggedRight,
  attach boxed title to top left={xshift=4mm,yshift=-2mm},
  boxed title style={arc=3mm, boxrule=0pt}
}

\newtcolorbox{infoband}[1]{
  enhanced,colback=#1!8,colframe=#1!50,arc=2mm,boxrule=.5pt,
  left=4mm,right=4mm,top=2mm,bottom=2mm,before skip=4pt,after skip=4pt,
  fontupper=\small\RaggedRight
}

% Puente narrativo entre secciones (contrato editorial Conéctate).
% Da continuidad: "Acabas de X. Ahora vas a Y porque Z."
% Visualmente: banda izquierda en color del grado, texto en itálica pequeña.
\newtcolorbox{puentebox}{
  enhanced,colback=milcGradeSoft,colframe=milcGradeDark,
  borderline west={3pt}{0pt}{milcGradeDark},
  arc=0mm,boxrule=0pt,left=5mm,right=4mm,top=2.5mm,bottom=2.5mm,
  before skip=4pt,after skip=8pt,
  fontupper=\itshape\small\color{milcNegro}\RaggedRight
}
\newcommand{\puente}[1]{%
  \begin{puentebox}%
    \textcolor{milcGradeDark}{\textbf{\textrightarrow}}\hspace{2mm}#1%
  \end{puentebox}%
}

\newcommand{\linead}[1]{\noindent\color{milcLinea}\rule{#1}{0.5pt}\color{milcNegro}}
\newcommand{\respuesta}[1]{\par\vspace{2mm}\foreach \n in {1,...,#1}{\noindent\color{milcLinea}\rule{\linewidth}{0.5pt}\par\vspace{5mm}}\color{milcNegro}}
\newcommand{\checkbox}{\textcolor{milcGradeDark}{\fbox{\phantom{x}}}\hspace{5pt}}

\newcommand{\phasecard}[4]{
\begin{tcolorbox}[enhanced,colback=#3,colframe=#2,arc=3mm,boxrule=.5pt,height=3.05cm,valign=center,halign=center,left=1.5mm,right=1.5mm,top=2mm,bottom=2mm]
{\sffamily\bfseries\fontsize{22}{24}\selectfont\textcolor{#2}{#1}}\par
{\sffamily\bfseries\small #4}
\end{tcolorbox}
}

\begin{document}
\thispagestyle{empty}
\begin{tikzpicture}[remember picture,overlay]
  \fill[white] (current page.south west) rectangle (current page.north east);
  \path[fill=milcGradePrimary,rounded corners=16mm]
    ([xshift=.55cm,yshift=-.55cm]current page.north west) rectangle
    ([xshift=-.55cm,yshift=3.65cm]current page.south east);

  \node[anchor=north west,text=milcGradeText] at ([xshift=2.0cm,yshift=-1.85cm]current page.north west)
    {\sffamily\bfseries\fontsize{14}{16}\selectfont GRADO};
  \node[anchor=north west,text=milcGradeText,opacity=.82] at ([xshift=2.0cm,yshift=-2.75cm]current page.north west)
    {\sffamily\bfseries\fontsize{9}{11}\selectfont SERIE GUÍAS · TECNOLOGÍA E INFORMÁTICA};

  \begin{scope}[shift={([xshift=-3.05cm,yshift=-2.55cm]current page.north east)}]
    \fill[white,opacity=.80] (-.62,-.52) rectangle (.62,.52);
    \draw[milcGradeText,line width=1pt] (-.62,-.52) rectangle (.62,.52);
    \fill[milcGradeDark] (-.42,-.38) rectangle (-.22,.06);
    \fill[milcGradeAccent] (-.10,-.38) rectangle (.10,.24);
    \fill[milcGradePrimary!60!black] (.22,-.38) rectangle (.42,.38);
  \end{scope}

  \node[anchor=west,text=milcGradeText] at ([xshift=1.9cm,yshift=-12.85cm]current page.north west)
    {\sffamily\bfseries\fontsize{124}{124}\selectfont 6};
  \draw[milcGradeText,line width=1.7pt]
    ([xshift=4.52cm,yshift=-11.68cm]current page.north west) circle (.13cm);

  \node[anchor=west,text=milcGradeText,text width=8.7cm,align=left] at ([xshift=10.35cm,yshift=-11.6cm]current page.north west)
    {
     {\sffamily\bfseries\fontsize{19}{22}\selectfont ¿Qué es\\internet? ---\\cómo viajan\\los datos\\a tu pantalla}\\[4mm]
     {\sffamily\bfseries\fontsize{23}{26}\selectfont Guía 27-6-TIC}\\[2mm]
     {\sffamily\fontsize{13}{16}\selectfont Período 3 · Procesador de texto e Internet}\\[5mm]
     {\sffamily\bfseries\fontsize{11}{13}\selectfont Institución Educativa Sor María Juliana}\\[1mm]
     {\sffamily\fontsize{10}{12}\selectfont Autor: Dr. Álvaro Cárdenas Orozco}
    };

  \path[fill=milcGradeSoft,rounded corners=7mm]
    ([xshift=1.35cm,yshift=.85cm]current page.south west) rectangle
    ([xshift=-1.35cm,yshift=4.25cm]current page.south east);
  \draw[milcGradeDark,line width=.65pt,rounded corners=7mm]
    ([xshift=1.35cm,yshift=.85cm]current page.south west) rectangle
    ([xshift=-1.35cm,yshift=4.25cm]current page.south east);
  \node[anchor=center,text=milcNegro,text width=17.0cm,align=center] at ([yshift=2.55cm]current page.south)
    {\sffamily\fontsize{10}{12}\selectfont Metodología MILC · Apertura ancestral · Pensamiento computacional · Triángulo Dussel-Estoicismo-Floridi\\
    \textcolor{milcGradeDark}{\bfseries Producto:} Un \textbf{diagrama del viaje de un paquete de datos} desde tu dedo cuando tocas un video de YouTube hasta el video aparece en tu pantalla. El diagrama incluye \textbf{6 estaciones} (tu dispositivo · módem/router de casa · proveedor de internet · cables submarinos · servidor de YouTube · regreso a tu pantalla), con flechas y tiempos aproximados. Más \textbf{en el cuaderno}: glosario de 6 términos clave (URL, IP, paquete, servidor, banda ancha, wifi) con definición tuya.};
\end{tikzpicture}
\mbox{}
\newpage

\section*{Apertura: ¿Qué es internet? --- cómo viajan los datos para llegar a tu pantalla}

\begin{softbox}{milcGradeDark}{milcGradeSoft}{Saber ancestral}
\textbf{Antes de internet, los mensajes viajaban por una red de personas y caballos.} El \textbf{pregonero} de Cartago llevaba un anuncio de cuadra en cuadra. Si la noticia tenía que llegar a Pereira, viajaba en \textbf{caballo} con un mensajero. Si tenía que llegar a Bogotá, viajaba en \textbf{tren}. Si llegaba a otro país, viajaba en \textbf{barco}. Cada mensajero conocía solo su tramo. \emph{Ningún mensajero hacía todo el viaje}: la información pasaba de uno a otro, como un balón. El cartero entregaba al jefe de correo, el jefe al maquinista, el maquinista al de la otra ciudad. \textbf{Esa idea --- la información que viaja en tramos, pasando de un sistema a otro --- es exactamente cómo funciona internet}. Solo que en vez de caballos hay cables; en vez de jinetes hay servidores; en vez de meses hay segundos. \textbf{Internet es el sistema de mensajeros más grande de la historia humana}, pero la idea es antigua: una red que conecta a personas que están lejos.
\end{softbox}

\begin{softbox}{milcTurquesa}{milcGris}{Saber contemporáneo}
\textbf{Internet} es una red mundial de computadoras conectadas entre sí. Cuando tú haces clic en un video de YouTube, NO le pides el video al video directamente. Le pides a tu \textbf{módem} (la cajita con luces en tu casa), el módem le pide al \textbf{proveedor} (Claro, Movistar, etc.), el proveedor le pide al \textbf{servidor de YouTube} (un computador grande en California o en otro país), y la respuesta hace el camino de vuelta. \textbf{La información viaja en pequeños trozos llamados ``paquetes''}, no de un solo bocado. Cada paquete tiene una \emph{dirección} (a quién va) y un \emph{remitente} (de dónde viene). Estas direcciones se llaman \textbf{IP} (Protocolo de Internet). Cada vez que abres una página, miles de paquetes viajan por cables submarinos, antenas, servidores. \textbf{Todo eso pasa en menos de 1 segundo}. Entender este viaje te quita el mito de que internet es \emph{``magia''}: es ingeniería bien hecha.
\end{softbox}

\begin{softbox}{milcMostaza}{milcGris}{Pregunta-puente}
\textit{Cuando le das \emph{play} a un video de YouTube, ¿\textbf{de dónde sale realmente el video}? ¿Está en tu celular? ¿Está en una nube en el cielo? ¿Dónde está físicamente?}
\end{softbox}

\begin{softbox}{milcVerde}{milcGris}{Saber y saber hacer}
Al terminar la clase: (1) podrás \textbf{identificar} las 6 estaciones del viaje de un paquete; (2) sabrás \textbf{explicar} qué son los paquetes, las direcciones IP y los servidores; (3) podrás \textbf{distinguir} entre wifi y datos móviles; (4) habrás \textbf{creado} un diagrama del viaje completo.
\end{softbox}

\small
\begin{tabularx}{\textwidth}{>{\bfseries}p{3.0cm}Y}
\rowcolor{milcGradePrimary!12}Autor & Dr. Álvaro Cárdenas Orozco\\
DBA & Produce contenidos digitales y valida información de fuentes web (MEN, T\&I 6°)\\
Referentes & Saber ancestral de la maestra rural · Lectura crítica · Soberanía informacional\\
Producto final & Un \textbf{diagrama del viaje de un paquete de datos} desde tu dedo cuando tocas un video de YouTube hasta el video aparece en tu pantalla. El diagrama incluye \textbf{6 estaciones} (tu dispositivo · módem/router de casa · proveedor de internet · cables submarinos · servidor de YouTube · regreso a tu pantalla), con flechas y tiempos aproximados. Más \textbf{en el cuaderno}: glosario de 6 términos clave (URL, IP, paquete, servidor, banda ancha, wifi) con definición tuya.\\
\end{tabularx}
\normalsize

\puente{Hoy haces 4 cosas: imaginas el viaje del paquete, aprendes las 6 estaciones, distingues conceptos clave (IP, paquete, servidor), y dibujas tu diagrama.}

\begin{titlebox}{milcGradeDark}
Ruta MILC: así vamos a aprender
\end{titlebox}
\begin{tabularx}{\textwidth}{>{\centering\arraybackslash}X>{\centering\arraybackslash}X>{\centering\arraybackslash}X>{\centering\arraybackslash}X}
\phasecard{1}{milcVerde}{milcGris}{Escuta\\[-1mm]\small Observo, escucho y pregunto.} &
\phasecard{2}{milcTurquesa}{milcGris}{Sistematización\\[-1mm]\small Pensamiento computacional.} &
\phasecard{3}{milcMagenta}{milcGris}{Praxis\\[-1mm]\small Construyo y aplico.} &
\phasecard{4}{milcVino}{milcGris}{Evaluación\\[-1mm]\small Reflexión liberadora.}
\end{tabularx}


\puente{Empezamos con una historia: el viaje de Mateo que envía un mensaje a su prima en España.}

\begin{titlebox}{milcVerde}
Fase 1 · Escuta
\end{titlebox}

\begin{softbox}{milcVerde}{milcGris}{Escena para escuchar}
\textbf{Actividad 1 · IDENTIFICA --- La historia del mensaje de Mateo} (10 min · individual). Lee esta historia y \textbf{en el cuaderno dibuja las estaciones que reconozcas}. \emph{``Mateo está en Cartago. Le quiere mandar una foto a su prima Sofía que vive en Madrid, España. Mateo abre WhatsApp en su celular. Toma la foto. La envía. En menos de 5 segundos, la foto aparece en el celular de Sofía. Pero la foto NO viajó directo de Cartago a Madrid. La foto pasó por: el wifi de su casa, el módem, el proveedor de internet de Cartago, cables que conectan Colombia con Europa por debajo del mar, otro proveedor en España, finalmente el celular de Sofía. ¿Cuántas estaciones intervinieron?''}. Cuéntalas. Dibújalas como un mapa.
\end{softbox}

\checkbox Leí la historia con calma.\\
\checkbox Conté las estaciones que mencionó.\\
\checkbox Dibujé un mini mapa con las estaciones que recuerdo.

\begin{infoband}{milcVerde}
\textbf{Lo que va al cuaderno (Actividad 1 · IDENTIFICA).} Título: «Actividad 1 --- El viaje de la foto de Mateo». \textbf{Formato:} dibujo simple de mapa con flechas conectando estaciones. \textbf{Extensión:} medio cuaderno. \textbf{Sabes que terminaste cuando} tu dibujo cubre el viaje de Cartago a Madrid con al menos 5 estaciones intermedias.
\end{infoband}

\puente{Después aprendes las 6 estaciones del viaje + los 6 términos técnicos.}

\begin{titlebox}{milcTurquesa}
Fase 2 · Sistematización con pensamiento computacional
\end{titlebox}

\begin{softbox}{milcTurquesa}{milcGris}{Cuatro pilares aplicados al tema}
Lo que dibujaste es básicamente correcto. Ahora aprende las \textbf{6 estaciones técnicas} con sus nombres reales y los términos que usan los ingenieros.
\end{softbox}

\small
\begin{tabularx}{\textwidth}{>{\bfseries\RaggedRight\arraybackslash}p{3.6cm}Y}
\rowcolor{milcTurquesa!90}\color{white}Pilar & \color{white}\bfseries Aplicación al tema\\
1. Descomposición & \textbf{Estación 1 --- Tu dispositivo (origen).} Tu celular, computador o tableta. Aquí nace el paquete: cuando tocas \emph{play}, tu dispositivo prepara una \textbf{solicitud} (\emph{request}) que dice: \emph{``quiero el video XYZ''}. La solicitud se divide en muchos \emph{paquetes} (pedacitos de información). \textbf{Estación 2 --- Módem o router de tu casa.} La cajita con luces. El módem recibe los paquetes de tu dispositivo (por wifi o cable) y los lanza hacia internet. Es como el \emph{cartero del barrio}: junta lo que sale de tu cuadra y lo manda al correo central.\\
2. Reconocimiento de patrones & \textbf{Estación 3 --- Proveedor de internet (ISP).} Claro, Movistar, ETB, etc. El proveedor recibe los paquetes del módem y los \textbf{enruta} hacia el destino correcto. Tiene grandes computadores (\emph{routers}) que deciden el mejor camino. Es como el \emph{centro de distribución de correo} de la ciudad. \textbf{Estación 4 --- Cables submarinos e infraestructura mundial.} Los paquetes viajan por \textbf{cables de fibra óptica} (algunos están en el fondo del mar) entre continentes. También por \emph{satélites} en ciertos casos. La velocidad: cerca de la velocidad de la luz. De Colombia a Europa, los paquetes tardan unos 100 milisegundos.\\
3. Abstracción & \textbf{Estación 5 --- Servidor de YouTube (destino).} \textbf{Un servidor es un computador grande} que guarda los videos y los entrega cuando alguien los pide. Los servidores de YouTube están en \emph{centros de datos} (\emph{data centers}) en California, Europa, Asia. Cada video tiene una \textbf{dirección IP} (un número que lo identifica único en internet, como una dirección postal). \textbf{Estación 6 --- El viaje de regreso.} El servidor manda los paquetes del video al revés: por los cables, por el proveedor de YouTube, por cables submarinos, por tu proveedor, por tu módem, hasta tu pantalla. Todo este viaje toma \textbf{menos de 1 segundo} para una página normal, 2-3 segundos para un video que empieza a cargar.\\
4. Algoritmo conceptual & \textbf{Los 6 términos técnicos clave.} (1) \textbf{URL}: la dirección que escribes en el navegador (\emph{www.youtube.com}). (2) \textbf{IP}: el número único de un dispositivo o servidor (\emph{172.217.16.78}). (3) \textbf{Paquete}: pedacito de información que viaja. Un video se divide en miles. (4) \textbf{Servidor}: computador grande que guarda y entrega información. (5) \textbf{Banda ancha}: cuánta información puede pasar por la conexión a la vez (más banda = más rápido). (6) \textbf{Wifi vs datos móviles}: wifi es de tu router de casa; datos móviles son del proveedor por antenas celulares. Ambos conectan a internet, solo cambia el primer tramo.\\
\end{tabularx}
\normalsize

\begin{drawbox}{milcTurquesa}{6 estaciones + 6 términos}
\textbf{Estaciones:} tu dispositivo · módem/router · proveedor · cables submarinos · servidor destino · regreso. \\[1mm]
\textbf{Términos:} URL · IP · paquete · servidor · banda ancha · wifi vs datos.
\end{drawbox}

\begin{softbox}{milcMostaza}{milcGris}{Errores comunes y su corrección}
\textbf{Mitos sobre internet que escucharás.} (1) \emph{``Los datos viajan por el aire''} --- Falso. Mayormente por cables (submarinos, terrestres). El aire (wifi/celular) es solo el último tramo. (2) \emph{``La nube está en el cielo''} --- Falso. La \emph{nube} son servidores en edificios físicos en distintos países. La palabra es metáfora. (3) \emph{``Cuanto más wifi, más internet''} --- Falso. Más wifi = más alcance dentro de tu casa, pero la velocidad la da el plan del proveedor. (4) \emph{``Internet es una sola cosa''} --- Falso. Es una red de redes (miles de proveedores y servidores conectados). De ahí su nombre: \emph{inter-net} (entre-red). (5) \emph{``Mi celular es el que guarda los videos de TikTok''} --- Falso. Los videos están en los servidores; tu celular solo descarga y reproduce.
\end{softbox}

\begin{infoband}{milcTurquesa}
\textbf{Lo que va al cuaderno (Actividad 2 · EXPLICA).} Título: «Actividad 2 --- Glosario de internet». \textbf{Formato:} lista de 6 términos con definición tuya en una línea. \textbf{Extensión:} 6 líneas. \textbf{Sabes que terminaste cuando} podrías explicarle a tu abuelo qué es una IP usando solo tus palabras.
\end{infoband}

\puente{Con todo claro, dibujas tu diagrama del viaje de un paquete específico (un video o un chat).}

\begin{titlebox}{milcMagenta}
Fase 3 · Praxis con pensamiento computacional
\end{titlebox}

\begin{softbox}{milcMagenta}{milcGris}{Construyendo el producto con los cuatro pilares}
\textbf{Actividad 3 · CREA --- Diagrama del viaje de un paquete} (28 min · individual). Vas a dibujar en una hoja del cuaderno el \textbf{viaje completo} de un paquete de datos desde que tocas \emph{play} en YouTube hasta que ves el video. 6 estaciones, flechas, tiempos.
\end{softbox}

\small
\begin{tabularx}{\textwidth}{>{\bfseries\RaggedRight\arraybackslash}p{3.6cm}Y}
\rowcolor{milcMagenta!90}\color{white}Pilar & \color{white}\bfseries Aplicación al producto\\
Descomposición & \textbf{Paso 1 --- Título y orientación.} En la parte superior escribe: \textbf{``El viaje de un paquete de datos''} y debajo: \emph{``Desde mi celular hasta el servidor de YouTube --- ida y vuelta''}. Subraya el título. Pon la hoja en \emph{horizontal} (apaisada) porque vas a necesitar espacio para 6 estaciones.\\
Patrones & \textbf{Paso 2 --- Las 6 estaciones (de izquierda a derecha).} Dibuja 6 cajitas o íconos en línea horizontal, separadas por flechas: \textbf{(1) Tu celular} (un cuadrito con ícono de móvil), \textbf{(2) Módem/router de casa} (una caja con antenas), \textbf{(3) Proveedor (Claro o Movistar)} (un edificio), \textbf{(4) Cables submarinos} (líneas onduladas que representan el océano), \textbf{(5) Servidor de YouTube} (un rectángulo grande con racks de computadores), \textbf{(6) Regreso} (puedes representarlo con una flecha curva grande que vuelve).\\
Abstracción & \textbf{Paso 3 --- Etiqueta cada estación.} Debajo de cada cajita, en tamaño pequeño, escribe: número (1, 2, 3...) + nombre + qué hace ahí en 1 línea. Ejemplo bajo la cajita 2: \emph{``Módem: junta los paquetes y los manda a internet''}.\\
Algoritmo de armado & \textbf{Paso 4 --- Tiempos aproximados.} En cada flecha entre estaciones, escribe el \textbf{tiempo aproximado} en milisegundos (ms). Datos reales aproximados: \emph{Celular → Módem: 1 ms · Módem → Proveedor: 10 ms · Proveedor → Cables submarinos: 30 ms · Cables → Servidor: 80 ms · Servidor procesa: 50 ms · Total ida: 170 ms · Total ida y vuelta: 340 ms}. (Total cerca de 1/3 de segundo).\\
\end{tabularx}
\normalsize

\begin{drawbox}{milcMagenta}{Antes de mostrar tu diagrama}
\checkbox El diagrama ocupa una hoja horizontal completa. \\[1mm]
\checkbox Están las 6 estaciones dibujadas y etiquetadas. \\[1mm]
\checkbox Hay flechas entre las 6 estaciones (ida y vuelta). \\[1mm]
\checkbox Cada flecha tiene un tiempo aproximado en ms. \\[1mm]
\checkbox Cada estación tiene una acción descrita en 1 línea. \\[1mm]
\checkbox Está firmado con frase de cierre subrayada.
\end{drawbox}

\begin{softbox}{milcMagenta}{milcGris}{Plantilla de trabajo}
\textbf{Estructura del diagrama.} \textit{Horizontal.} Estaciones 1 a 5 en línea. Flechas con ms. Etiquetas con número + nombre + acción. Estación 6 (regreso) como flecha curva. Firma y frase abajo.
\end{softbox}

\begin{infoband}{milcMagenta}
\textbf{Lo que va al cuaderno (Actividad 3 · CREA).} Título: «Actividad 3 --- El viaje de un paquete de datos». \textbf{Formato:} diagrama horizontal con 6 estaciones + flechas + tiempos. \textbf{Extensión:} 1 hoja completa horizontal. \textbf{Sabes que terminaste cuando} podrías mostrarle el diagrama a tu mamá y ella entendería cómo funciona internet en 2 minutos.
\end{infoband}

\puente{El producto es ese diagrama firmado + el glosario en el cuaderno.}

\begin{titlebox}{milcMostaza}
Producto de la sesión
\end{titlebox}
\textbf{Diagrama del viaje de un paquete con 6 estaciones + tiempos · firmado}.
\begin{softbox}{milcMostaza}{milcGris}{Criterios de calidad}
(1) El diagrama tiene 6 estaciones identificables. (2) Las flechas muestran el orden del viaje (ida y vuelta). (3) Cada estación tiene número, nombre y acción. (4) Hay tiempos aproximados en milisegundos. (5) Está firmado con frase de cierre.
\end{softbox}

\puente{Antes de salir, verificas que tu diagrama tiene 6 estaciones y que cada flecha tiene un tiempo aproximado.}

\begin{titlebox}{milcVino}
Fase 4 · Evaluación liberadora — cinco dimensiones
\end{titlebox}

\begin{softbox}{milcVino}{milcGris}{Sentido de la evaluación}
Evaluar no es solo revisar una nota. Es reconocer qué aprendí, qué cambió en mí, qué emociones manejé, cómo me relaciono con la ciudad y la comunidad, y qué le dejaré a quien viene detrás.
\end{softbox}

\textbf{Mi reflexión liberadora en cinco dimensiones.} Completa cada frase iniciada:

\small
\begin{tabularx}{\textwidth}{>{\bfseries\RaggedRight\arraybackslash}p{3.4cm}Y>{\RaggedRight\arraybackslash}p{4.6cm}}
\rowcolor{milcVino}\color{white}Dimensión & \color{white}\bfseries Mi respuesta (frame starter) & \color{white}\bfseries Algo que descubrí\\
1. Personal & Lo que más cambió en mí fue \linead{6.5cm} & \linead{4.2cm}\\
2. Emocional & La emoción más fuerte fue \linead{2.5cm} y la regulé \linead{3cm} & \linead{4.2cm}\\
3. Ciudadana & En democracia esto importa porque \linead{6cm} & \linead{4.2cm}\\
4. Local (Cartago) & En mi barrio sería útil para \linead{6.5cm} & \linead{4.2cm}\\
5. Intergeneracional & A mi abuela/abuelo le diría \linead{6.5cm} & \linead{4.2cm}\\
\end{tabularx}
\normalsize

\begin{infoband}{milcVino}
\textbf{Para la próxima sustentación.} Vuelve a esta tabla en seis meses. Verás que las respuestas cambian — no porque lo aprendido cambie, sino porque tú cambias. La evaluación liberadora es una conversación contigo a través del tiempo.
\end{infoband}


\puente{Cerramos con 3 voces que te ayudan a entender por qué saber cómo funciona la red te hace ciudadano digital.}

\begin{titlebox}{milcGradeDark}
Triángulo de pensamiento · cierre filosófico
\end{titlebox}

\begin{softbox}{milcVino}{milcGris}{Dussel · lente del nosotros}
\textit{````La tecnología se vuelve servidumbre cuando no la entiendes. Se vuelve libertad cuando sabes cómo funciona.''''}

Mucha gente usa internet sin saber qué pasa por dentro. Eso los hace \textbf{usuarios dependientes}: si algo falla, no saben qué reportar; si los engañan, no saben cómo. \emph{Cuando tú entiendes el viaje del paquete, dejas de ser usuario y pasas a ser ciudadano de la red}. La diferencia es enorme. Hoy diste un paso grande hacia esa ciudadanía.

\textbf{¿Quiero seguir siendo usuario dependiente o quiero ser ciudadano de la red?}
\end{softbox}
\linead{\linewidth}\\[1mm]
\linead{\linewidth}

\begin{softbox}{milcMostaza}{milcGris}{Estoicismo · Marco Aurelio (emperador que se preguntaba por la naturaleza de las cosas) · cuidado interior}
\textit{````Cada cosa cotidiana esconde un mundo si te detienes a preguntarla. Internet es un mundo entero detrás de un clic.''''}

Hacer clic en \emph{play} es un gesto que toma 1 segundo. Pero \textbf{detrás de ese clic} hay cables submarinos, servidores en California, ingenieros que diseñaron los protocolos hace 50 años, electricidad que mueve los routers, infraestructura por todo el planeta. \emph{Detenerse a entenderlo es entrenar el asombro}. Cuando vuelves a usar internet con esa conciencia, todo cambia.

\textbf{¿Qué otras cosas cotidianas escondo mundos que no he querido conocer?}
\end{softbox}
\linead{\linewidth}\\[1mm]
\linead{\linewidth}

\begin{softbox}{milcTurquesa}{milcGris}{Floridi · lente de la infoesfera}
\textit{````Internet no es una herramienta más. Es el sistema nervioso de la humanidad del siglo XXI. Aprender cómo funciona es aprender en qué planeta vives.''''}

\textbf{La generación de tus papás vivió un mundo sin internet}. La tuya nació con él. Eso te da una ventaja: lo usas con naturalidad. Pero también un riesgo: \emph{lo das por sentado}. Entender cómo funciona te ayuda a no darlo por sentado: a notar cuando falla, a defender que sea libre y abierto, a usarlo con criterio. Internet es público, como el agua, como la energía. \emph{Entenderlo es responsabilidad ciudadana}.

\textbf{¿Estoy usando internet como un usuario de paso, o como un ciudadano que sabe lo que tiene entre manos?}
\end{softbox}
\linead{\linewidth}\\[1mm]
\linead{\linewidth}

\begin{titlebox}{milcVino}
Compromiso final
\end{titlebox}

\small
\begin{tabularx}{\textwidth}{>{\RaggedRight\arraybackslash}p{3.0cm}YYY}
\rowcolor{milcVino}\color{white}\bfseries Criterio & \color{white}\bfseries Lo logré & \color{white}\bfseries En proceso & \color{white}\bfseries Necesito apoyo\\
Comprensión del tema & Explico ideas centrales con mis palabras. & Reconozco ideas pero falta precisión. & Necesito repasar conceptos base.\\
Pensamiento computacional & Apliqué los 4 pilares al diseñar mi producto. & Apliqué algunos pilares. & Necesito releer la fase 2.\\
Aplicación y producto & Mi producto cumple los criterios y se puede revisar. & Producto funcional con ajustes pendientes. & Aún en construcción.\\
Reflexión 5 dimensiones & Respondí cada dimensión con honestidad. & Respondí algunas con detalle. & Mis respuestas son muy generales.\\
Triángulo de pensamiento & Conecté las 3 voces con mi trabajo real. & Conecté al menos una. & No relaciono las voces con mi proceso.\\
\end{tabularx}
\normalsize

\textbf{Hoy entendí que:}\\[1mm]
\linead{\linewidth}\\[1mm]
\linead{\linewidth}

\textbf{Mi compromiso para la próxima sesión es:}\\[1mm]
\linead{\linewidth}\\[1mm]
\linead{\linewidth}

\begin{softbox}{milcMostaza}{milcGris}{Cita de despedida · Marco Aurelio}
\textit{``El obstáculo en el camino se vuelve el camino. Lo que estorba al camino se vuelve camino.''}\\[1mm]
Cada error de hoy, bien observado, se convierte en pieza del camino que pisarás mañana.
\end{softbox}

\small
\begin{tabularx}{\textwidth}{>{\bfseries}p{3.6cm}Y}
\rowcolor{milcGradeDark!12}Nombre completo: & \linead{10cm}\\
Fecha: & \linead{4cm} \quad Curso/grupo: \linead{4cm}\\
Mi firma: & \linead{9cm}\\
\end{tabularx}
\normalsize

\begin{infoband}{milcGradeDark}
\textbf{Para la próxima sesión.} Esta semana, observa: cuando navegas, ¿cuánto se demora cada página en cargar? ¿Por qué algunas tardan más? La próxima clase aprendes a \textbf{buscar bien en internet}: operadores de búsqueda, filtros, técnicas para encontrar lo que necesitas en segundos, no en minutos.
\end{infoband}

\end{document}
